% !TeX root = ../main.tex
\section{Kế hoạch Marketing và Chiến lược thâm nhập thị trường}

% !TeX root = ../main.tex
\subsection{Chiến lược tiếp thị và Kênh truyền thông}

\subsubsection{Thông điệp truyền thông cốt lõi}

Mutux được định vị là nền tảng giúp người dùng \textbf{trải nghiệm thiết bị gaming chuyên sâu trước khi quyết định sở hữu}, đồng thời giúp chủ sở hữu khai thác hiệu quả thiết bị nhàn rỗi. Thông điệp này được triển khai nhất quán cùng các cơ chế tạo niềm tin gồm xác minh danh tính (KYC), mức ký quỹ phù hợp với giá trị tài sản và quy trình ghi nhận tình trạng thiết bị minh bạch. Nhờ đó, hoạt động truyền thông không chỉ giới thiệu lợi ích của việc thuê mà còn giải đáp trực tiếp rào cản lớn nhất của mô hình chia sẻ tài sản: niềm tin giữa các bên giao dịch.

\begin{itemize}
    \item \textbf{Slogan chủ đạo: ``Gear Up. Play On.''} Khẩu hiệu ngắn gọn, giàu tính hành động và phù hợp với văn hóa gaming/esports. Thông điệp bổ trợ \textit{``Thuê thử gear chuẩn gu, mua khi đã chắc''} được sử dụng tại các điểm chạm hướng đến người thuê để nhấn mạnh lợi ích giảm rủi ro khi lựa chọn thiết bị.

    \item \textbf{Giá trị cho người thuê (Renter): ``Trải nghiệm chuyên sâu, lựa chọn có chủ đích.''} Người thuê có thể tiếp cận thiết bị ở phân khúc cao với chi phí ban đầu hợp lý, sử dụng trong điều kiện thực tế và đưa ra quyết định mua phù hợp hơn. Cơ chế \textit{Rent-to-Own} --- nếu được triển khai sau giai đoạn kiểm chứng --- cho phép khấu trừ một phần chi phí thuê vào giá mua theo chính sách công bố trước; đây là một lựa chọn bổ sung, không phải điều kiện bắt buộc của giao dịch thuê.

    \item \textbf{Giá trị cho người cho thuê (Lender): ``Bảo vệ tài sản, khai thác giá trị.''} Với cá nhân, Mutux hỗ trợ biến thiết bị ít sử dụng thành nguồn thu bổ sung trong một quy trình có kiểm soát rủi ro. Với SME, Mutux có thể trở thành kênh ``trải nghiệm trước khi mua'', cung cấp dữ liệu phản hồi thực tế để hỗ trợ bán hàng. Trọng tâm truyền thông cho phía cung là sự an toàn, minh bạch và hiệu quả kinh tế.
\end{itemize}

\subsubsection{Phễu chuyển đổi và Phân bổ nguồn lực theo giai đoạn}

Chiến lược tiếp thị được thiết kế theo phễu chuyển đổi gồm: \textbf{nhận biết} (nội dung và cộng đồng), \textbf{cân nhắc} (trang thiết bị, đánh giá và giải đáp chính sách), \textbf{kích hoạt} (đăng ký, KYC và đơn thuê đầu tiên), \textbf{giữ chân} (trải nghiệm tốt, ưu đãi tái thuê) và \textbf{giới thiệu} (đánh giá, liên kết chia sẻ). Mỗi kênh đều gắn mã định danh chiến dịch và một lời kêu gọi hành động cụ thể, giúp Mutux đo được đường đi từ lượt tiếp cận đến giao dịch hoàn tất.

Trong sáu tháng đầu, ngân sách và nhân sự tiếp thị được phân bổ theo nguyên tắc ``kiểm chứng trước, mở rộng sau'' như sau:

\begin{longtable}{@{}p{0.17\textwidth}p{0.23\textwidth}p{0.20\textwidth}p{0.28\textwidth}@{}}
\toprule
\textbf{Giai đoạn} & \textbf{Tỷ trọng nguồn lực} & \textbf{Kênh trọng tâm} & \textbf{Mục tiêu và điều kiện chuyển pha} \\
\midrule
\endhead
Tháng 1--2: kiểm chứng & 70\% & Cộng đồng trường đại học, Discord, tiếp cận trực tiếp và nội dung thử nghiệm ngắn & Xây dựng nhóm người dùng lõi, kiểm chứng nhu cầu và quy trình KYC--bàn giao. Chỉ chuyển pha khi có phản hồi tích cực từ cả người thuê lẫn người cho thuê và quy trình vận hành ổn định. \\
Tháng 3--4: tối ưu & 20\% & TikTok, Facebook và tái tiếp cận người đã quan tâm & So sánh hiệu quả thông điệp, tối ưu trang thiết bị và ưu đãi đơn đầu. Chỉ tăng chi khi CAC, tỷ lệ hoàn tất KYC và tỷ lệ đơn hoàn tất đạt ngưỡng nội bộ đã phê duyệt. \\
Tháng 5--6: mở rộng có chọn lọc & 10\% & Nhân rộng kênh hiệu quả, referral và hợp tác SME thử nghiệm & Tăng nguồn cung và giao dịch mà vẫn duy trì chất lượng dịch vụ, tỷ lệ hủy và tỷ lệ tranh chấp trong giới hạn chấp nhận được. \\
\bottomrule
\end{longtable}

\paragraph{Kênh tiếp cận và giữ chân người thuê (Renter)}

\begin{itemize}
    \item \textbf{TikTok --- ``Trải nghiệm thực tế, đánh giá chân thực'':} Phát triển các chuỗi ``So sánh thông số trong sử dụng thực tế'', ``Bộ gear theo nhu cầu'' và video bàn giao minh họa quy trình an toàn. Mỗi nội dung dẫn đến một trang thiết bị hoặc landing page có mã chiến dịch riêng; lời kêu gọi hành động là xem tình trạng thiết bị, đăng ký hoặc hoàn tất KYC.

    \item \textbf{Facebook và cộng đồng trường --- bằng chứng xã hội và giải đáp chính sách:} Facebook đóng vai trò điểm chạm để công khai đánh giá, giải đáp về KYC, ký quỹ và xử lý sự cố. Các câu lạc bộ gaming/công nghệ tại trường được tiếp cận qua hoạt động dùng thử quy mô nhỏ, nhằm tạo bằng chứng xã hội từ trải nghiệm thật thay vì chỉ phụ thuộc vào quảng cáo trả phí.

    \item \textbf{Mutux Guild trên Discord --- cộng đồng lõi:} Discord được sử dụng để duy trì tương tác sâu với nhóm người dùng am hiểu gaming, không phải kênh tiếp cận đại chúng duy nhất. Tại khu vực thí điểm như ĐHQG-HCM, thành viên Guild có thể tham gia buổi trải nghiệm, hướng dẫn kiểm tra thiết bị và đóng góp phản hồi cho sản phẩm. Vòng lặp mong muốn là: trải nghiệm tốt $\rightarrow$ đánh giá xác thực $\rightarrow$ đơn thuê mới.

    \item \textbf{Giữ chân và điểm tín nhiệm:} Sau khi kiểm chứng tính khả thi về tài chính và quản trị rủi ro, Mutux có thể thử nghiệm điểm tín nhiệm cho người thuê có lịch sử thanh toán, bàn giao và bảo quản tốt. Điểm này có thể được quy đổi thành ưu đãi tái thuê hoặc giảm \textit{mức ký quỹ yêu cầu} trong phạm vi chính sách rủi ro, thay vì giảm ký quỹ một cách tự động.
\end{itemize}

\paragraph{Kênh tiếp cận nguồn cung (Lender \& SME)}

\begin{itemize}
    \item \textbf{Tiếp cận trực tiếp chủ sở hữu và SME:} Tập trung vào chủ sở hữu thiết bị chất lượng, câu lạc bộ công nghệ và cửa hàng/SME phù hợp với khu vực thí điểm. Thông điệp chào mời tập trung vào quy trình bảo vệ tài sản, khả năng tiếp cận khách hàng mới và bảng tổng hợp hiệu quả danh mục.

    \item \textbf{Gói ``Thuê thử cùng Mutux'' cho SME:} Đây là gói hợp tác thử nghiệm, trong đó đối tác đưa một số thiết bị vào nền tảng để người dùng trải nghiệm trước khi mua. Báo cáo gửi đối tác gồm lượt xem, lượt thuê, tỷ lệ chuyển đổi thuê sang mua (nếu có) và phản hồi tổng hợp; mọi việc chia sẻ dữ liệu phải tuân thủ sự đồng ý của người dùng và chính sách bảo mật.

    \item \textbf{Referral cho Lender:} Người cho thuê có thể chia sẻ liên kết danh mục thiết bị. Phần thưởng giới thiệu chỉ được áp dụng cho giao dịch hợp lệ, hoàn tất và không phát sinh tranh chấp trong thời hạn quy định. Cách tiếp cận này biến nguồn cung chất lượng thành một kênh giới thiệu có thể đo lường, đồng thời kiểm soát chi phí thu hút.

    \item \textbf{``Community Test Vault'' --- thử nghiệm có kiểm soát:} Chương trình được thí điểm với số lượng thiết bị giới hạn từ chủ sở hữu đã xác minh. Mỗi thiết bị có hồ sơ tình trạng công khai, quy định bàn giao và trách nhiệm rõ ràng. Mục tiêu là tạo không gian trải nghiệm đáng tin cậy, đồng thời kiểm chứng liệu mô hình này có giúp nâng chất lượng nguồn cung hay không.
\end{itemize}

\subsubsection{Quản trị rủi ro truyền thông}

Vì Mutux xử lý giao dịch tài sản có giá trị và ký quỹ, phản hồi truyền thông phải dựa trên bằng chứng và bảo vệ công bằng cho cả hai bên. Mọi tranh chấp về tình trạng thiết bị hoặc hoàn trả ký quỹ được xử lý theo nguyên tắc nhất quán, có lưu vết và được thông báo rõ cho các bên liên quan.

\begin{itemize}
    \item \textbf{Phòng ngừa:} Công bố trước giao dịch các điều khoản ký quỹ, tiêu chuẩn ghi nhận tình trạng thiết bị, quy trình khiếu nại và thời hạn xử lý. Nội dung hướng dẫn được chuẩn hóa cho cả người thuê và người cho thuê.
    \item \textbf{Tiếp nhận và xác minh:} Phản hồi ban đầu trong vòng 24 giờ làm việc; thu thập ảnh, video hoặc biên bản bàn giao; đối chiếu với điều khoản đã công bố. Thời gian phản hồi ban đầu không đồng nghĩa với thời hạn kết luận tranh chấp.
    \item \textbf{Cập nhật và khắc phục:} Cập nhật trạng thái cho các bên theo mốc thời gian xác định, chỉ công bố thông tin đã được xác minh và bảo vệ dữ liệu cá nhân. Khi xảy ra sự cố có ảnh hưởng rộng, Mutux đưa ra thông báo nhất quán về sự việc, biện pháp khắc phục và cách liên hệ hỗ trợ nhằm duy trì niềm tin của cộng đồng.
\end{itemize}

\subsubsection{Hệ thống chỉ số đo lường Hiệu quả (KPIs)}

Mutux ưu tiên chỉ số phản ánh giao dịch, thanh khoản marketplace và chất lượng dịch vụ thay cho các chỉ số bề nổi. Các chỉ số được theo dõi theo kênh, khu vực thí điểm và nhóm người dùng để xác định nguyên nhân thay vì chỉ quan sát giá trị trung bình.

\begin{longtable}{@{}p{0.22\textwidth}p{0.43\textwidth}p{0.23\textwidth}@{}}
\toprule
\textbf{Nhóm KPI} & \textbf{Định nghĩa và phương pháp đo lường} & \textbf{Mục tiêu quản trị} \\
\midrule
\endhead
Phễu người thuê & Tỷ lệ từ lượt truy cập $\rightarrow$ đăng ký $\rightarrow$ hoàn tất KYC $\rightarrow$ đơn thuê hoàn tất, theo từng mã chiến dịch. & Xác định điểm thắt và đánh giá hiệu quả từng kênh. \\
Thanh khoản marketplace & Số listing hoạt động, tỷ lệ listing có đơn, thời gian để lender có đơn đầu tiên và số đơn hoàn tất trên mỗi thiết bị hoạt động. & Bảo đảm phát triển cân bằng giữa cung và cầu. \\
Hiệu quả thu hút & CAC = tổng chi phí tiếp thị phân bổ / số người thuê mới có đơn hoàn tất; theo dõi thêm tỷ lệ hoàn vốn CAC và LTV/CAC khi đủ dữ liệu. & Chỉ mở rộng ngân sách ở kênh tạo tăng trưởng bền vững. \\
Giữ chân và giới thiệu & Tỷ lệ tái thuê trong các mốc thời gian xác định, tỷ lệ người dùng đánh giá sau giao dịch và tỷ lệ đơn hợp lệ qua referral. & Đo giá trị dài hạn và sức mạnh truyền miệng. \\
Chất lượng dịch vụ và rủi ro & Tỷ lệ hủy, tỷ lệ tranh chấp, thời gian phản hồi ban đầu, tỷ lệ hoàn tất bàn giao có bằng chứng. & Bảo vệ trải nghiệm và uy tín của nền tảng khi tăng trưởng. \\
\bottomrule
\end{longtable}

Trong giai đoạn kiểm chứng, mức chuyển đổi từ đăng ký sang đơn thuê hoàn tất $5\%$--$10\%$ được xem là \textbf{giả định để kiểm chứng}, không phải cam kết doanh thu. Ngưỡng cụ thể cho CAC, biên lợi nhuận ròng trên mỗi giao dịch, tỷ lệ hủy và tỷ lệ tranh chấp sẽ được thiết lập cùng mô hình tài chính và vận hành trước khi tăng ngân sách. Quyết định mở rộng chỉ được đưa ra khi đồng thời thỏa mãn các điều kiện về hiệu quả thu hút, thanh khoản marketplace và chất lượng dịch vụ.


% !TeX root = ../main.tex
\subsection{Kế hoạch ra mắt và chiến lược phát triển thị trường}

\subsubsection{Chiến lược tiếp cận người dùng đầu tiên}

Vì Mutux là marketplace kết nối người thuê và người cho thuê, nhóm không nên tiếp cận toàn bộ thị trường ngay từ đầu. Chiến lược phù hợp hơn là chọn một cụm nhỏ quanh ĐHQG-HCM, ký túc xá, CLB eSports và một số shop gaming gear lân cận. Khi trong cùng một khu vực đã có cả nguồn cung và người thuê đầu tiên, Mutux mới có cơ sở mở rộng sang các cộng đồng khác \cite{rochetTirole2003,chen2021coldstart}.

Trong giai đoạn đầu, Mutux tập trung vào ba nhóm gear dễ giải thích và dễ trải nghiệm: chuột gaming, bàn phím cơ và tai nghe. Nhóm sáng lập có thể chủ động mời một số shop SME và người sưu tầm gear tham gia ``Founding Shelf'', sau đó tuyển người thuê qua CLB, nhóm bạn, ký túc xá và các cộng đồng gaming. Việc hỗ trợ trực tiếp cách đăng thiết bị, chụp ảnh và kể câu chuyện món gear phù hợp với bài học khởi động thủ công của Airbnb \cite{graham2013scale}.

\subsubsection*{Câu chuyện marketing trung tâm: ``Một chiếc gear, hai câu chuyện''}

Mutux có thể kể câu chuyện về một món gear đang nằm yên ở kệ shop hoặc trong bộ sưu tập cá nhân, trong khi một sinh viên đang cần thử thiết bị trước một giải đấu, một dự án hoặc quyết định mua hàng. Mutux giúp hai nhân vật gặp nhau: người cho thuê có thêm cơ hội khai thác tài sản, còn người thuê có trải nghiệm thực tế để biết thiết bị có phù hợp hay không.

Một nội dung mẫu có thể theo chân Minh Quân: Quân xem nhiều review nhưng vẫn chưa biết con chuột mới có hợp tay mình. Sau một buổi ``Gear Match Day'', Quân thuê thiết bị trong vài ngày, ghi lại điều mình thích và không thích, rồi chia sẻ kết luận cho người chơi tiếp theo. Như vậy, câu chuyện không cần kết thúc bằng việc ``gear giúp thắng'', mà bằng việc người dùng đưa ra lựa chọn tự tin hơn. Cách tiếp cận này phù hợp với nghiên cứu về storytelling, trải nghiệm thương hiệu và truyền miệng \cite{escalas2004narrative,brakus2009brandexperience,nielsen2021trust}.

\subsubsection{Kế hoạch phát hành và ra mắt sản phẩm}

Do sản phẩm vẫn ở giai đoạn MVP, Mutux nên ra mắt theo hình thức thí điểm nhỏ trước khi mở rộng. Kế hoạch có thể gồm bốn bước:

\begin{longtable}{@{}p{0.16\textwidth}p{0.25\textwidth}p{0.40\textwidth}@{}}
\toprule
\textbf{Giai đoạn} & \textbf{Mục tiêu} & \textbf{Hoạt động chính} \\
\midrule
\endhead
Chuẩn bị & Có nguồn cung và câu chuyện đầu tiên & Mời một số shop/collector tham gia, chọn những mẫu gear phù hợp và mở danh sách người quan tâm. \\
Thí điểm kín & Kiểm chứng trải nghiệm thật & Mời nhóm nhỏ người thuê dùng thử trong vài ngày; nhóm sáng lập hỗ trợ trực tiếp và thu nhận xét trung thực. \\
Ra mắt trong trường & Tạo điểm gặp giữa hai phía & Tổ chức ``Gear Match Day'' tại CLB, phòng máy hoặc không gian sinh viên; kết hợp dùng thử, kể chuyện chủ gear và đăng ký thuê. \\
Tiếp tục và mở rộng & Duy trì truyền miệng & Đăng các nhật ký trải nghiệm, Gear Drop theo tuần và chương trình giới thiệu; chỉ mở sang khu vực mới khi cụm đầu tiên vận hành ổn định. \\
\bottomrule
\end{longtable}

Việc hợp tác với CLB và giải đấu sinh viên là hướng tiếp cận phù hợp vì đây là nơi tập trung sẵn người chơi, đội tuyển và khán giả. Các hoạt động như USEC Fresher và HCMUS Cup cho thấy cộng đồng eSports sinh viên đã có những điểm kết nối cụ thể tại khu vực Mutux lựa chọn \cite{vngUsecFresher2025,vngHcmusCup2026}. Mutux nên ưu tiên tài trợ trải nghiệm hoặc gear dùng thử thay vì chỉ đặt logo.

\subsubsection{Hoạt động tiếp thị giai đoạn đầu}

Các hoạt động dưới đây được xem là những cách triển khai khác nhau của cùng một câu chuyện. Nhóm không cần thực hiện đồng thời tất cả, mà có thể chọn một số hoạt động phù hợp với nguồn lực của môn học.

\begin{longtable}{@{}p{0.22\textwidth}p{0.49\textwidth}p{0.16\textwidth}@{}}
\toprule
\textbf{Ý tưởng} & \textbf{Cách triển khai} & \textbf{Mục tiêu} \\
\midrule
\endhead
``Gear ngủ quên'' & Mời shop và collector kể về những món gear ít được sử dụng; Mutux biến các câu chuyện này thành nội dung giới thiệu nguồn cung. & Thu hút lender \\
``72 giờ tìm đúng gear'' & Một sinh viên ghi lại trải nghiệm trước, trong và sau khi thuê; nội dung được khuyến khích trung thực cả khi thiết bị không phù hợp. & Tạo review thật \\
Gear Match Day & Buổi dùng thử tại trường/CLB, có khu thử mù, tường chia sẻ chuyện mua nhầm và bàn tư vấn từ người đã dùng. & Tạo trải nghiệm \\
Gear Blind Date & Người tham gia thử chuột, bàn phím hoặc tai nghe trước khi biết thương hiệu và giá. & Tạo tò mò \\
Gear Captain & Mỗi CLB có một thành viên làm đầu mối mời người dùng, tổ chức demo và thu phản hồi. & Tiếp cận cộng đồng \\
Demo Shelf cùng shop & Một góc ``Thử thật cùng Mutux'' tại shop, kết hợp lịch thuê thử và chính sách mua sau thuê nếu đối tác có công bố rõ ràng. & Tạo nguồn cung neo \\
Pass the Play & Sau một giao dịch hoàn tất, người thuê mời đồng đội tiếp theo; ưu đãi chỉ phát sinh khi đơn hợp lệ hoàn thành. & Tạo truyền miệng \\
Tournament Emergency Kit & Phối hợp với CLB/ban tổ chức để giới thiệu một số gear dự phòng trước giải, không cam kết luôn có thiết bị thay thế. & Gắn với nhu cầu cấp thiết \\
Gear Passport & Một thẻ giấy ghi câu chuyện của món gear; mỗi người thuê có thể thêm một nhận xét ngắn cho người kế tiếp. & Tạo cảm giác tiếp nối \\
Developer Desk Week & Ở giai đoạn sau, thử nghiệm chuột công thái học, bàn phím yên tĩnh hoặc tai nghe với lập trình viên trẻ. & Mở rộng phân khúc \\
\bottomrule
\end{longtable}

Ngoài các hoạt động trên, Mutux có thể thử thêm ``Gear Drop thứ Sáu'', ``Squad Test Pack'' cho nhóm bạn, ``Gear sai gu confession'', hoặc ``Lender Open House''. Đây là các biến thể nhỏ của cùng một hướng marketing: để người dùng thật kể câu chuyện về thiết bị thay cho việc chỉ quảng cáo thông số. Các chương trình nội dung cộng đồng như GoPro Awards và mô hình campus ambassador của Red Bull có thể được dùng làm nguồn tham khảo, nhưng không nên sao chép quy mô của họ \cite{goproAwards2015,redbullStudentMarketeer}.

Trong phạm vi môn học, mục tiêu của kế hoạch là chứng minh Mutux biết chọn thị trường ban đầu, xây dựng câu chuyện phù hợp và đề xuất cách kiểm chứng nhu cầu. Các chỉ tiêu về số người dùng, chi phí marketing hay doanh thu chỉ nên được xem là giả định để thảo luận, không phải cam kết triển khai. Mutux cũng không nên hứa ``gear giúp thắng'', ``thu nhập bảo đảm'' hoặc ``không cần cọc''; niềm tin cần được xây dựng từ trải nghiệm và đánh giá trung thực \cite{godesMayzlin2009,ryuFeick2007referral}.


